\begin{table}[tbp]
\centering
\caption{问题三：五个监测感官维度的劣化}
\label{tab:t09_q3_defects}
\footnotesize
\setlength{\tabcolsep}{3.5pt}
\begin{tabular}{llrrr}
\toprule
方法 & 感官维度 & 劣化 n/N & 劣化率 & Δ 中位数 \\
\midrule
无校正 & 老化酒花 & 527/924 & 57.03\% & 0.172 \\
无校正 & DMS & 462/924 & 50.00\% & 0.020 \\
无校正 & 双乙酰 & 449/924 & 48.59\% & -0.100 \\
无校正 & 马厩味 & 554/924 & 59.96\% & 0.178 \\
无校正 & 金属味 & 561/924 & 60.71\% & 0.127 \\
无校正 & 任一维 & 868/924 & 93.94\% &  \\
OT 投影 & 老化酒花 & 519/924 & 56.17\% & 0.171 \\
OT 投影 & DMS & 455/924 & 49.24\% & 0.013 \\
OT 投影 & 双乙酰 & 457/924 & 49.46\% & -0.079 \\
OT 投影 & 马厩味 & 551/924 & 59.63\% & 0.176 \\
OT 投影 & 金属味 & 559/924 & 60.50\% & 0.124 \\
OT 投影 & 任一维 & 869/924 & 94.05\% &  \\
联合 QP & 老化酒花 & 0/924 & 0.00\% & -0.233 \\
联合 QP & DMS & 0/924 & 0.00\% & -0.740 \\
联合 QP & 双乙酰 & 0/924 & 0.00\% & -0.068 \\
联合 QP & 马厩味 & 0/924 & 0.00\% & -0.234 \\
联合 QP & 金属味 & 0/924 & 0.00\% & -0.100 \\
联合 QP & 任一维 & 0/924 & 0.00\% &  \\
\bottomrule
\end{tabular}
\end{table}
% 注（写入图注/表注，勿直接编译）：Δ = 约束后 − 约束前的风格内标准化感官分；Δ > 0 表示该缺陷维度变差。五个维度是预注册的监测维度，未使用测试集或生成结果反向挑选。
